\section{Matrículas especiais}

\subsection{Veículos importados (1997 e 1998)}\label{sec:importados}

Chapas de formato e cores normais com a sequência \plate{00-00-KA} a
\plate{00-00-KF} foram reservadas aos veículos importados em segunda mão
matriculados durante 1997.

Nos veículos importados por particulares entre 1997 e 1998, a primeira letra do
grupo de letras também era um \plate{K}; a combinação mais alta atribuída terá
sido \plate{KG}.\footnote{A Wikipédia indica que todos os grupos \plate{KA} a \plate{KZ}
foram atribuídos a importações em 1997; as duas fontes não coincidem.}

\platefig{import-8949kc}{Veículo importado: \plate{89-49-KC}.}

\subsection{Veículos novos ainda não matriculados}

Chapas de fundo vermelho com caracteres a branco, sempre com 220~mm $\times$
330~mm, para veículos novos ainda não matriculados. Na parte de cima têm um
número sequencial e na de baixo a denominação (ou abreviatura) do operador
registado do veículo, do fabricante ou do concessionário.

\platefig[0.3\linewidth]{new-6019-santagri}{Veículo novo ainda não matriculado.}

\subsection{Veículos de competição}

Chapas de fundo vermelho com caracteres a preto, parecidas com as das máquinas e
atrelados industriais, mas com o ano e o mês da primeira matrícula.

\platefig{competition-39th48}{Veículo de competição: \plate{39-TH-48}.}

\subsection{Veículos para exportação}

Chapas de fundo amarelo com caracteres a preto, com a sequência \plate{00000-X},
em que \plate{X} identifica a alfândega onde o veículo foi registado. Na barra
amarela, por cima dos números da data, são acrescentadas as letras \plate{EXP}.

\gen{export}

\platefig[0.35\linewidth]{export-00573l}{Veículo para exportação: \plate{00573-L}.}

\subsection{Máquinas e atrelados industriais}

Chapas de fundo vermelho com caracteres a preto. Na barra amarela, em vez dos
números da data, é colocada uma letra que indica a classe de circulação da
máquina ou do atrelado (classes \plate{A} a \plate{H}). A atribuição destas matrículas é regulada pelo Decreto-Lei n.º~107/2006 (RAMMI). As chapas têm a sequência normal dos automóveis.

\platepair{industrial-gp1234}{Máquina industrial: \plate{GP-12-34}, classe \plate{A}.}%
          {industrial-48xj68}{Máquina industrial: \plate{48-XJ-68}, classe \plate{A}.}

\subsection{Tratores agrícolas}

Os tratores agrícolas e os seus reboques têm chapas de fundo verde (Decreto
Regulamentar n.º~13/98).

\subsection{Guarda Nacional Republicana}

Chapas de formato e cores normais com a sequência \plate{GNR X-0} a
\plate{GNR X-0000}, em que \plate{X} é uma letra que indica o tipo ou a função do
veículo (ver o anexo~\ref{sec:gnr}).

\subsection{Forças militares}

\begin{description}
  \item[Exército Português] Chapas de formato e cores normais com as sequências
    \plate{MG-00-00}, \plate{ME-00-00} e \plate{MX-00-00}. A maioria ainda usa as
    chapas pretas com letras brancas.
  \item[Marinha Portuguesa] Sequência \plate{AP-00-00}.
  \item[Força Aérea Portuguesa] Sequência \plate{AM-00-00}.
  \item[Forças aliadas] Configuração semelhante à das matrículas normais, mas com
    a letra \plate{U} na eurobanda em vez do \plate{P}. São usadas sobretudo por
    militares norte-americanos na Base das Lajes, na ilha Terceira, desde 2005.
\end{description}

\platepair{army-mx5531}{Exército: \plate{MX-55-31}.}%
          {allied-b589an}{Forças aliadas: \plate{B5-89-AN}, com \plate{U} na eurobanda.}

\subsection{Veículos do Presidente da República}

São talvez as chapas mais distintas das especiais e são usadas em todos os carros
afetos ao Presidente da República: fundo preto com o escudo nacional ao centro,
ladeado pelas letras \plate{P} e \plate{R} a branco ou prateado.

\platefig[0.3\linewidth]{president}{Chapa dos veículos do Presidente da República.}

\subsection{Corpo diplomático}\label{sec:diplomaticos}

Placas brancas com caracteres a vermelho, sem eurobanda, com a sequência
\plate{000-XX000}, reservadas a diplomatas e a funcionários em missões
internacionais. O Decreto-Lei n.º~54/2005 (art.~6.º) estende os caracteres, traços
e rebordo vermelhos às entidades abrangidas pelo Protocolo relativo aos Privilégios
e Imunidades das Comunidades Europeias. O primeiro grupo de algarismos identifica o país ou a
organização internacional (ver o anexo~\ref{sec:codigos-diplomaticos}) e
\plate{XX} pode ser:
\begin{itemize}
  \item \plate{CD}: Corpo Diplomático;
  \item \plate{CC}: Corpo Consular;
  \item \plate{FM}: Funcionário em Missão Internacional.
\end{itemize}

\platepair{diplomatic-cd}{Corpo Diplomático: \plate{206-CD444}.}%
          {diplomatic-cc}{Corpo Consular: \plate{024-CC410}.}
